\chapter{Ficando com as Migalhas}
\label{cap:migalhas}

Finalização do rito endereçando o que sobra, pois algo sempre sobra, tanto a se perguntar, quanto a se recusar;

Nego Bispo falando do que fazer com o cocô de galinha;
    o que apodrece se joga no mato para virar adubo, e o que não estraga se guarda:

Fragmentos também são migalhas, como o "kit-de-odores" o foi para adubar o totem;

\textbf{-x-} a contraparte, o contraponto, o contratempo -- complemento -- negativo como positivo

\textbf{-x-} as cartas que restaram no jogo sem se utilizar -- anti-biblioteca de Umberto Eco

---

Os contrapontos, as incompletudes, o espaço negativo como espaço positivo -- articular é ver o que falta e o que sobra.

As migalhas são as contingências ainda não nomeadas;
    aquilo que Latour rejeita por não articular bem:
        o rito começa a terminar endereçando as migalhas enquanto o que sobre e o que falta na realização do rito.

O leitor acompanha numa dinâmica com as cartas;
    composição de uma anti-biblioteca:
        o que apodrece se joga no mato para virar adubo, o que não estraga se guarda.

Ética articulada;
    não é sobre identificar o que é verdade:
        é sobre saber se deslocar-localizar entre sistemas de verdade distintos -- é sobre saber aterrar.
